%% Software and Systems Modeling (SoSyM) manuscript
%% Based on the Springer Nature LaTeX template, version 3.1 (December 2024),
%% kept unmodified in template/. sn-jnl.cls and sn-mathphys-num.bst sit next to
%% this file because the submission must include all style files.
%%
%% SoSyM submission guidelines (checked 2026-09-13):
%%   - LaTeX with the [iicol] option; upload the full source and a compiled PDF
%%   - abstract of 150 to 250 words, without undefined abbreviations or citations
%%   - numbered citations in square brackets (sn-mathphys-num)
%%   - decimal headings, at most three levels
%%   - a "Statements and Declarations" section with Competing Interests,
%%     plus a data availability statement
%%   - single-blind review, so author names and repository links can stay
%% The template asks for a single .tex file, so avoid \input.

\documentclass[iicol,pdflatex,sn-mathphys-num]{sn-jnl}

%%%% Standard packages from the template
\usepackage{graphicx}%
\usepackage{multirow}%
\usepackage{amsmath,amssymb,amsfonts}%
\usepackage{amsthm}%
\usepackage{mathrsfs}%
\usepackage[title]{appendix}%
\usepackage{xcolor}%
\usepackage{textcomp}%
\usepackage{manyfoot}%
\usepackage{booktabs}%
\usepackage{algorithm}%
\usepackage{algorithmicx}%
\usepackage{algpseudocode}%
\usepackage{listings}%

%% Code listings for Java and XMI snippets
\lstset{basicstyle=\ttfamily\footnotesize,columns=fullflexible,
  breaklines=true,frame=single,captionpos=b}

\theoremstyle{thmstylethree}%
\newtheorem{definition}{Definition}%

%% Drafting aid: remove every \todo before submission
\newcommand{\todo}[1]{\textcolor{red}{[TODO: #1]}}

\raggedbottom

\begin{document}

\title[Decentralized Model Persistence]{Towards Decentralized Model Persistence}

\author*[1]{\fnm{Alfa} \sur{Yohannis}}\email{alfa.ryano@pradita.ac.id}
%% Further authors: \author[1]{\fnm{Given} \sur{Family}}\email{name@example.org}

\affil*[1]{\orgdiv{Department}, \orgname{Pradita University}, \orgaddress{\city{City}, \country{Indonesia}}}

\abstract{\todo{150 to 250 words: the problem of location-based model
persistence, the content-addressed alternative, the answers to RQ1--RQ3, and
the main quantitative result.}}

\keywords{Model-driven engineering, Model persistence, Eclipse Modeling Framework, IPFS, Content-addressed storage, Cross-resource references}

\maketitle

\section{Introduction}\label{sec:introduction}

\todo{Motivation, problem statement, contributions, and paper outline. The
introduction should not have subheadings.}

%% Research questions, verbatim from research_questions.md.
%% conversation_summary.md suggests RQ2 for MODELS and RQ1/RQ3 for the
%% journal version; drop or shorten RQ2 here if it is published separately.
\begin{enumerate}
\item[RQ1] What are the architectural challenges of extending model persistence frameworks to support content-addressed decentralized storage, and how can they be addressed?
\item[RQ2] How can cross-resource references be maintained and resolved in an immutable, content-addressed environment where resource identifiers change on every modification?
\item[RQ3] What is the performance overhead of content-addressed model persistence, and under what conditions is it practical for MDE toolchains?
\end{enumerate}

\section{Background}\label{sec:background}

\subsection{Model persistence in EMF}\label{sec:background-emf}

\todo{Resource, ResourceSet, URIConverter, proxies and lazy resolution in
EMF~\cite{steinberg2008emf}; location-based URIs such as file:// and
platform://.}

\subsection{Content-addressed storage with IPFS}\label{sec:background-ipfs}

\todo{CIDs, immutability, IPNS names, and the Kubo HTTP API~\cite{benet2014ipfs}.}

\section{Content-addressed model persistence}\label{sec:rq1}

%% RQ1. Sources: context.md, core-plugin/
\subsection{Architecture}\label{sec:rq1-architecture}

\todo{IPFSResourceImpl, IPFSResourceFactoryImpl and IPFSURIHandlerImpl, and
how ipfs:// URIs plug into URIConverter and the resource factory registry.}

\subsection{Resource lifecycle}\label{sec:rq1-lifecycle}

\todo{Create, save (new CID), modify, save again (another CID).}

\subsection{Mixing file-based and content-addressed resources}\label{sec:rq1-mixed}

\todo{One ResourceSet holding both XMIResource and IPFSResource instances.}

\section{Cross-resource references}\label{sec:rq2}

%% RQ2. Sources: research_questions.md, ipfs_advantages.md
\subsection{Version pinning by content addressing}\label{sec:rq2-pinning}

\todo{Why a CID reference pins the exact version of the referenced resource.}

\subsection{Reference management strategies}\label{sec:rq2-strategies}

\todo{Direct CID references, IPFS directories, IPNS indirection, and a
manifest or registry, with their trade-offs.}

\subsection{Cascade save}\label{sec:rq2-cascade}

\todo{The cascadeSave() algorithm, cycle protection, and the stale-CID case for
bidirectional references.}

\section{Tool support in the Eclipse BPMN2 Modeler}\label{sec:tool}

%% Sources: README.md of org.eclipse.bpmn2-modeler-1.5.4-ecorefs
\todo{Opening and publishing models and projects from the BPMN2 IPFS menu,
the JSON project manifest, IPNS publishing, and per-project Kubo settings.}

\section{Evaluation}\label{sec:evaluation}

\subsection{Setup}\label{sec:evaluation-setup}

\todo{Kubo daemon in Docker, Java runtime, hardware, warm-up and measured
iterations.}

\subsection{RQ1: Functional validation}\label{sec:evaluation-rq1}

\todo{Save and load round trips and proxy resolution over ipfs:// URIs.}

\subsection{RQ2: Reference evolution}\label{sec:evaluation-rq2}

\todo{Version-pin test, intentional upgrade test, cascadeSave latency, and the
BPMN case study with cross-file references.}

\subsection{RQ3: Performance}\label{sec:evaluation-rq3}

%% Results: eval/results/
\todo{Save and load latency against local XMI files for models of increasing
size, with mean, median, and standard deviation.}

\subsection{Large-scale models: MoDisco models of Eclipse JDT}\label{sec:evaluation-modisco}

%% Sources: eval/modisco_jdt_reverse/README.md and MANUAL-CHECKS.md
\todo{Partitioning by compilation unit, component manifests published under
IPNS, and EMF Compare checks against snapshots fetched from IPFS.}

\section{Discussion}\label{sec:discussion}

\todo{Limitations: no lazy loading, no query support, IPNS resolution
latency, cascade cost, and circular references.}

\subsection{Threats to validity}\label{sec:threats}

\todo{Internal, external, and construct validity.}

\section{Related work}\label{sec:related}

\todo{XMI file persistence, CDO, NeoEMF, and model versioning approaches.}

\section{Conclusion}\label{sec:conclusion}

\todo{Answers to RQ1--RQ3 and future work.}

\backmatter

\bmhead{Acknowledgements}

\todo{Optional.}

\section*{Statements and Declarations}

\bmhead{Funding}
\todo{Funding statement.}

\bmhead{Competing interests}
\todo{Competing interests statement.}

\bmhead{Ethics approval and consent to participate}
\todo{Write ``Not applicable'' if it does not apply.}

\bmhead{Consent for publication}
\todo{Write ``Not applicable'' if it does not apply.}

\bmhead{Data availability}
\todo{Where the benchmark results and model datasets can be obtained.}

\bmhead{Code availability}
The implementation is available at
\url{https://github.com/alfa-yohannis/ecorefs}, the BPMN2 Modeler integration at
\url{https://github.com/alfa-yohannis/org.eclipse.bpmn2-modeler-1.5.4-ecorefs},
and the BPMN~2.0 metamodel changes at
\url{https://github.com/alfa-yohannis/org.eclipse.bpmn2-ecorefs}.

\bmhead{Author contribution}
\todo{Contribution of each author.}

\bibliography{references}

\end{document}
